\chapter{Antiphospholipid Syndrome}
\index{Antiphospholipid Syndrome}

\section*{Definition and Overview}
Antiphospholipid syndrome (APS) is an \textbf{autoimmune disorder} in which the immune system produces antibodies that mistakenly attack proteins that bind to phospholipids -- normal components of cell membranes and blood components. These antibodies make the blood abnormally prone to \textbf{clotting}, causing blockages in veins and arteries, and they also cause complications in pregnancy, including recurrent miscarriage and pre-eclampsia.

APS is also known as \textbf{Hughes syndrome}. It occurs either on its own (\textbf{primary APS}) or together with another autoimmune disease, most commonly systemic lupus erythematosus (\textbf{secondary APS}). A rare and severe variant, \textbf{catastrophic APS}, presents with widespread clotting in small blood vessels and multiple organ failure. APS is an important, treatable cause of stroke in younger people and of recurrent pregnancy loss.

\section*{Epidemiology}
APS is uncommon but not rare; estimates suggest that between 1 in 500 and 1 in 1000 people are affected. It occurs more often in women than in men, and its prevalence is substantially higher in people with lupus, in whom the condition coexists in roughly one-third of cases. The typical age of presentation is in young and middle adulthood.

Antiphospholipid antibodies can be detected transiently in healthy people during infections and in some drug exposures without ever causing clinical problems, which is why laboratory findings alone are not enough for a diagnosis. APS is responsible for a significant minority of strokes and transient ischaemic attacks in people under 50, and it is identified in a substantial proportion of women investigated for recurrent pregnancy loss.

\section*{Aetiology and Pathophysiology}
In APS, antibodies are directed against proteins that bind phospholipids, particularly \textbf{beta-2-glycoprotein I}. The exact sequence by which these antibodies cause clotting is incompletely understood, but they are thought to activate the endothelial cells lining blood vessels, prime platelets, and interfere with natural anticoagulant proteins, tipping the balance of the clotting cascade toward clot formation.

The consequences depend on where clotting occurs:

\begin{itemize}
    \item \textbf{Venous clots:} most often deep vein thrombosis of the legs, which can travel to the lungs as a pulmonary embolism
    \item \textbf{Arterial clots:} most seriously in the brain, causing stroke or transient ischaemic attacks; they can also affect the heart, kidneys, and limbs
    \item \textbf{Pregnancy:} clots in the placental circulation can cause recurrent miscarriage, stillbirth, pre-eclampsia, and restricted fetal growth
    \item \textbf{Microvascular clotting:} in catastrophic APS, widespread small-vessel occlusion causes multiple organ failure
    \item \textbf{Triggers:} infection, surgery, immobility, pregnancy, and oestrogen-containing contraceptives can provoke the first or a subsequent event
\end{itemize}

\section*{Clinical Features}
\begin{itemize}
    \item \textbf{Deep vein thrombosis:} swelling, pain, warmth, and redness of one calf or thigh, often after immobility
    \item \textbf{Pulmonary embolism:} sudden breathlessness, pleuritic chest pain, coughing up blood, or collapse
    \item \textbf{Stroke or transient ischaemic attack:} sudden weakness or numbness, facial droop, speech disturbance, or visual loss, occurring at a younger age than is typical for stroke
    \item \textbf{Pregnancy complications:} recurrent early miscarriage, late pregnancy loss, stillbirth, severe pre-eclampsia, or poor fetal growth
    \item \textbf{Livedo reticularis:} a mottled, bluish, lace-like discolouration of the skin, especially on the arms and legs
    \item \textbf{Other features:} migraine-like headaches, low platelet count (thrombocytopenia), leg ulcers, superficial thrombophlebitis, and heart valve thickening
\end{itemize}

\section*{Differential Diagnosis}
\begin{itemize}
    \item \textbf{Inherited thrombophilias:} factor V Leiden, prothrombin gene mutation, and protein C, protein S, or antithrombin deficiency
    \item \textbf{Acquired causes of clotting:} cancer, surgery, immobility, pregnancy, and oestrogen-containing medication without underlying APS
    \item \textbf{Other causes of stroke in the young:} cardiac emboli, carotid or vertebral dissection, and migraine-related ischaemia
    \item \textbf{Other autoimmune diseases:} lupus, vasculitis, and other systemic inflammatory conditions
    \item \textbf{Other causes of recurrent pregnancy loss:} chromosomal abnormalities, uterine structural problems, and endocrine disorders
    \item \textbf{Transient positive antibodies:} present during infection without clinical significance
\end{itemize}

\section*{When to Seek Medical Care}
Anyone with known APS who becomes pregnant, or who is planning a pregnancy, should be referred early to a specialist obstetric team, because treatment during pregnancy differs from routine management.

\textbf{Contact your local emergency services for:}
\begin{itemize}
    \item Sudden breathlessness, chest pain, or coughing up blood -- possible pulmonary embolism
    \item Sudden weakness or numbness of a limb, facial droop, or difficulty speaking -- possible stroke
    \item A swollen, painful, red, and warm leg, especially after immobility -- possible deep vein thrombosis
    \item Fainting, collapse, or a very fast or irregular heartbeat
    \item High fever with confusion, widespread clotting problems, or organ failure -- possible catastrophic APS
\end{itemize}

\section*{Diagnosis and Investigations}
The diagnosis of APS requires both clinical features and persistent positive antibodies. A single positive blood test is not sufficient.

\begin{itemize}
    \item \textbf{Antibody tests:} the lupus anticoagulant, anticardiolipin antibodies, and anti-beta-2-glycoprotein I antibodies
    \item \textbf{Repeat testing:} positive antibodies must be confirmed on two occasions at least 12 weeks apart, with the clinical event, to satisfy diagnostic criteria
    \item \textbf{Imaging:} venous ultrasound for DVT; CT pulmonary angiography or ventilation-perfusion scanning for pulmonary embolism; CT or MRI of the brain for stroke
    \item \textbf{Full blood count:} to detect thrombocytopenia and haemolytic anaemia, which can accompany the syndrome
    \item \textbf{Coagulation screen:} with attention to the effect of anticoagulants, which can interfere with lupus anticoagulant testing
    \item \textbf{Specialist review:} confirmation and ongoing management by a haematologist or rheumatologist
\end{itemize}

\section*{Management}
\begin{itemize}
    \item \textbf{Anticoagulation after a clot:} treatment begins with heparin and continues with warfarin, which reduces the risk of further clots markedly; novel oral anticoagulants are used selectively in some cases under specialist guidance
    \item \textbf{Lifelong treatment:} after a first clot, most people remain on anticoagulation long term because recurrence risk is high if it is stopped
    \item \textbf{Pregnancy:} low-dose aspirin and low-molecular-weight heparin injections under specialist care, because warfarin is avoided in pregnancy
    \item \textbf{Avoid oestrogen:} combined oral contraceptives and hormone replacement therapy are generally avoided because they increase clotting risk
    \item \textbf{Control other risk factors:} blood pressure, cholesterol, diabetes, and smoking all add to the risk of arterial and venous clots
    \item \textbf{Compression stockings and mobility:} reduce the risk of leg clots during long journeys or immobility
    \item \textbf{Multidisciplinary care:} haematologists, rheumatologists, and obstetricians working together, with regular monitoring of anticoagulation intensity
\end{itemize}

\section*{Complications}
\begin{itemize}
    \item Recurrent blood clots, most commonly in the months after a first event or if anticoagulation is stopped
    \item Pulmonary embolism, which can be fatal
    \item Stroke with permanent neurological disability
    \item Recurrent pregnancy loss, stillbirth, pre-eclampsia, and fetal growth restriction
    \item Chronic kidney disease and pulmonary hypertension from repeated small-vessel clots
    \item Heart valve disease and accelerated atherosclerosis
    \item Catastrophic APS, a rare emergency characterised by multi-organ failure
    \item Bleeding complications from anticoagulant treatment itself
\end{itemize}

\section*{Prognosis and Outcomes}
With proper anticoagulation, most people with APS live full, active lives and the risk of further clots is greatly reduced. The highest risk of recurrence is soon after a first event and whenever treatment is interrupted. Women who have had thrombotic APS are usually treated with anticoagulation throughout pregnancy and the postnatal period, but with aspirin and low-molecular-weight heparin the majority of pregnancies end successfully.

Catastrophic APS remains a life-threatening emergency despite aggressive treatment with anticoagulation, steroids, and plasma exchange. Long-term monitoring is needed because APS is a lifelong condition, but with adherence to treatment and control of other cardiovascular risk factors, most people achieve good outcomes.

\section*{Prevention and Self-Care}
\begin{itemize}
    \item Take anticoagulant medicines exactly as prescribed and never stop them without medical advice
    \item Avoid smoking and keep blood pressure, cholesterol, and weight controlled
    \item Avoid oestrogen-containing contraceptives and hormone replacement therapy
    \item Move regularly during long journeys and consider compression stockings on flights or during bed rest
    \item Seek early, specialist-led care if pregnant or planning a pregnancy
    \item Carry a medical alert card so emergency staff know about the condition and its treatment
    \item Report symptoms of a clot (leg pain, breathlessness) or of unusual bleeding promptly
    \item Attend scheduled anticoagulation monitoring and follow-up appointments
\end{itemize}

\section*{Sources and Further Reading}
The following organisations provide reliable, up-to-date information on antiphospholipid syndrome for patients, students, and health professionals.

\begin{itemize}
    \item NHS. \textit{Information on antiphospholipid syndrome and its management.} https://www.nhs.uk/conditions/
    \item Mayo Clinic. \textit{Overview of antiphospholipid syndrome, symptoms, and treatment.} https://www.mayoclinic.org/
    \item MedlinePlus. \textit{Health information on antiphospholipid syndrome and blood clotting.} https://medlineplus.gov/
    \item MSD Manual. \textit{Professional overview of the antiphospholipid syndrome.} https://www.msdmanuals.com/
    \item NICE. \textit{Guidance relevant to venous thromboembolism and anticoagulation.} https://www.nice.org.uk/
    \item World Health Organization. \textit{Resources on chronic and autoimmune diseases.} https://www.who.int/
\end{itemize}
